% Appendix: proof of Theorem~\ref{thm:stable_sparse_hybrid_residual}
% following Candes, Romberg, and Tao, arXiv:math/0503066, Section~2.2 (Theorem~2).
% Block size is \mathsf{M}=3|T_0| to avoid confusion with the measurement matrix $M$.

\section{Complete proof of Theorem~\ref{thm:stable_sparse_hybrid_residual}}
\label{sec:appendix_bpdn_rip_proof}

Fix $A\in\mathbb{R}^{m\times N}$, $\alpha^\star\in\mathbb{R}^N$, and $z=A\alpha^\star+w$ with $\|w\|_2\le \varepsilon$.
Let $\widehat{\alpha}$ solve \eqref{eq:bpdn} with the same $\varepsilon$, and set $h:=\widehat{\alpha}-\alpha^\star$.
Let $T_0\subset\{1,\ldots,N\}$ index $k$ largest magnitudes of $\alpha^\star$ (ties broken arbitrarily), so $|T_0|=k$.
Let $\alpha^\star_k$ agree with $\alpha^\star$ on $T_0$ and vanish on $T_0^c$; this is a best $k$-sparse approximation in the usual basis model.
Then $\|\alpha^\star_{T_0^c}\|_1=\|\alpha^\star-\alpha^\star_k\|_1$.
Define the block size
\begin{equation}
  \mathsf{M} := 3|T_0|
  \label{eq:def_block_size_Msf}
\end{equation}
and $\rho := |T_0|/\mathsf{M}=1/3$.
Partition $T_0^c$ into disjoint sets $T_1,T_2,\ldots$ of cardinalities $\mathsf{M}$ (except possibly the last), arranged so that $(|h_i|)_{i\in T_1}$ dominates $(|h_i|)_{i\in T_2}$, etc., on $T_0^c$.
Set $T_{01}:=T_0\cup T_1$.

\subsection{Tube and cone constraints}

\begin{lemma}[Tube constraint]
\label{lem:appendix_tube}
One has $\|Ah\|_2\le 2\varepsilon$.
\end{lemma}

\begin{proof}
By the triangle inequality,
$\|Ah\|_2=\|A(\widehat{\alpha}-\alpha^\star)\|_2\le \|A\widehat{\alpha}-z\|_2+\|z-A\alpha^\star\|_2\le 2\varepsilon$.
\end{proof}

\begin{lemma}[Cone constraint]
\label{lem:appendix_cone_compressible}
One has
\begin{equation}
  \|h_{T_0^c}\|_1 \le \|h_{T_0}\|_1 + 2\|\alpha^\star_{T_0^c}\|_1.
  \label{eq:appendix_cone_compressible}
\end{equation}
\end{lemma}

\begin{proof}
Feasibility of $\alpha^\star$ and optimality of $\widehat{\alpha}$ for the $\ell_1$ objective give $\|\widehat{\alpha}\|_1\le \|\alpha^\star\|_1$.
Write $\widehat{\alpha}=\alpha^\star+h$.
For any index set $T$,
\[
  \|\alpha^\star+h\|_1
  =
  \|(\alpha^\star+h)_T\|_1 + \|(\alpha^\star+h)_{T^c}\|_1
  \ge
  \|\alpha^\star_T\|_1 - \|h_T\|_1 + \|h_{T^c}\|_1 - \|\alpha^\star_{T^c}\|_1,
\]
by the reverse triangle inequality on each block.
Thus
\[
  \|\alpha^\star\|_1
  \ge
  \|\widehat{\alpha}\|_1
  \ge
  \|\alpha^\star_{T_0}\|_1 - \|h_{T_0}\|_1 + \|h_{T_0^c}\|_1 - \|\alpha^\star_{T_0^c}\|_1.
\]
Rearranging and using $\|\alpha^\star\|_1=\|\alpha^\star_{T_0}\|_1+\|\alpha^\star_{T_0^c}\|_1$ yields \eqref{eq:appendix_cone_compressible}.
\end{proof}

\subsection{Tail bounds}

\begin{lemma}[Summed tail block norms]
\label{lem:appendix_sum_tail}
One has
\begin{equation}
  \sum_{j\ge 2}\|h_{T_j}\|_2 \le \frac{\|h_{T_0^c}\|_1}{\sqrt{\mathsf{M}}}.
  \label{eq:appendix_sum_tail_blocks}
\end{equation}
\end{lemma}

\begin{proof}
For $j\ge 2$, each $p\in T_j$ satisfies $|h_p|\le \min_{q\in T_{j-1}}|h_q|\le \|h_{T_{j-1}}\|_1/\mathsf{M}$.
Hence $\|h_{T_j}\|_2^2\le |T_j|\cdot (\|h_{T_{j-1}}\|_1/\mathsf{M})^2\le \|h_{T_{j-1}}\|_1^2/\mathsf{M}$, so $\|h_{T_j}\|_2\le \|h_{T_{j-1}}\|_1/\sqrt{\mathsf{M}}$.
Summing over $j\ge 2$ yields
\[
  \sum_{j\ge 2}\|h_{T_j}\|_2
  \le
  \frac{1}{\sqrt{\mathsf{M}}}\sum_{j\ge 1}\|h_{T_j}\|_1
  =
  \frac{\|h_{T_0^c}\|_1}{\sqrt{\mathsf{M}}}.
\]
\end{proof}

\begin{lemma}[Tail $\ell_2$ norm on $T_{01}^c$]
\label{lem:appendix_tail_T01c}
Define $\eta:=\|\alpha^\star_{T_0^c}\|_1/\sqrt{|T_0|}$.
Then
\begin{equation}
  \|h_{T_{01}^c}\|_2 \le \frac{\|h_{T_0}\|_1 + 2\|\alpha^\star_{T_0^c}\|_1}{\sqrt{\mathsf{M}}}
  \le
  \sqrt{\rho}\,\bigl(\|h_{T_0}\|_2 + 2\eta\bigr).
  \label{eq:appendix_tail_T01c_bound}
\end{equation}
\end{lemma}

\begin{proof}
Since $T_{01}^c=\bigcup_{j\ge 2}T_j$, one has $\|h_{T_{01}^c}\|_2\le \sum_{j\ge 2}\|h_{T_j}\|_2$.
Apply Lemma~\ref{lem:appendix_sum_tail} and Lemma~\ref{lem:appendix_cone_compressible} to get
$\|h_{T_{01}^c}\|_2\le (\|h_{T_0}\|_1+2\|\alpha^\star_{T_0^c}\|_1)/\sqrt{\mathsf{M}}$.
Finally, $\|h_{T_0}\|_1\le \sqrt{|T_0|}\,\|h_{T_0}\|_2$ and $\|\alpha^\star_{T_0^c}\|_1=\sqrt{|T_0|}\,\eta$.
\end{proof}

\begin{lemma}[Global $\ell_2$ bound]
\label{lem:appendix_global_l2}
One has
\begin{equation}
  \|h\|_2 \le (1+\sqrt{\rho})\,\|h_{T_{01}}\|_2 + 2\sqrt{\rho}\,\eta.
  \label{eq:appendix_global_l2}
\end{equation}
\end{lemma}

\begin{proof}
By the triangle inequality, $\|h\|_2\le \|h_{T_{01}}\|_2+\|h_{T_{01}^c}\|_2$.
Also $\|h_{T_0}\|_2\le \|h_{T_{01}}\|_2$.
Insert Lemma~\ref{lem:appendix_tail_T01c}.
\end{proof}

\subsection{Sensing lower bound on the head block}

\begin{lemma}[Head sensing inequality]
\label{lem:appendix_head_sensing}
Assume RIP of orders $\mathsf{M}+|T_0|$ and $\mathsf{M}$ with constants $\delta_{\mathsf{M}+|T_0|}$ and $\delta_{\mathsf{M}}$.
Define
\begin{equation}
  C_{|T_0|,\mathsf{M}}
  :=
  \sqrt{1-\delta_{\mathsf{M}+|T_0|}}
  -
  \sqrt{\rho}\,\sqrt{1+\delta_{\mathsf{M}}}.
  \label{eq:def_C_TM}
\end{equation}
Then
\begin{equation}
  \|Ah\|_2
  \ge
  C_{|T_0|,\mathsf{M}}\,\|h_{T_{01}}\|_2
  -
  2\sqrt{\rho}\,\sqrt{1+\delta_{\mathsf{M}}}\,\eta.
  \label{eq:appendix_Ah_lower}
\end{equation}
\end{lemma}

\begin{proof}
Since $h=h_{T_{01}}+\sum_{j\ge 2}h_{T_j}$, the reverse triangle inequality gives
\[
  \|Ah\|_2
  \ge
  \|A h_{T_{01}}\|_2 - \sum_{j\ge 2}\|A h_{T_j}\|_2.
\]
The vector $h_{T_{01}}$ is $(\mathsf{M}+|T_0|)$-sparse and each $h_{T_j}$ is $\mathsf{M}$-sparse, so RIP implies
$\|A h_{T_{01}}\|_2\ge \sqrt{1-\delta_{\mathsf{M}+|T_0|}}\,\|h_{T_{01}}\|_2$
and $\|A h_{T_j}\|_2\le \sqrt{1+\delta_{\mathsf{M}}}\,\|h_{T_j}\|_2$.
Lemma~\ref{lem:appendix_sum_tail} and Lemma~\ref{lem:appendix_cone_compressible} yield
\[
  \sum_{j\ge 2}\|h_{T_j}\|_2
  \le
  \frac{\|h_{T_0}\|_1+2\|\alpha^\star_{T_0^c}\|_1}{\sqrt{\mathsf{M}}}
  \le
  \sqrt{\rho}\,\bigl(\|h_{T_0}\|_2+2\eta\bigr),
\]
as in the proof of Lemma~\ref{lem:appendix_tail_T01c}.
Combine the estimates.
\end{proof}

\subsection{Positivity and conclusion}

\begin{lemma}[Positivity of $C_{|T_0|,\mathsf{M}}$]
\label{lem:appendix_C_positive}
If $\delta_{3k}+3\delta_{4k}<2$ and $|T_0|=k$, then $C_{|T_0|,\mathsf{M}}>0$.
\end{lemma}

\begin{proof}
Here $\mathsf{M}=3|T_0|=3k$, so $\mathsf{M}+|T_0|=4k$ and $\delta_{\mathsf{M}+|T_0|}=\delta_{4k}$, $\delta_{\mathsf{M}}=\delta_{3k}$.
Thus
\[
  C_{|T_0|,\mathsf{M}}
  =
  \sqrt{1-\delta_{4k}} - \sqrt{\tfrac{1}{3}}\,\sqrt{1+\delta_{3k}}.
\]
Squaring the strict inequality $\sqrt{1-\delta_{4k}}>\sqrt{(1+\delta_{3k})/3}$ (both sides nonnegative) and rearranging yields $2>\delta_{3k}+3\delta_{4k}$, which is assumed.
\end{proof}

\begin{proof}[Proof of Theorem~\ref{thm:stable_sparse_hybrid_residual}]
Lemma~\ref{lem:appendix_tube} gives $\|Ah\|_2\le 2\varepsilon$.
Lemma~\ref{lem:appendix_head_sensing} combined with Lemma~\ref{lem:appendix_C_positive} implies
\[
  \|h_{T_{01}}\|_2
  \le
  \frac{2}{C_{|T_0|,\mathsf{M}}}\,\varepsilon
  +
  \frac{2\sqrt{\rho}\,\sqrt{1+\delta_{\mathsf{M}}}}{C_{|T_0|,\mathsf{M}}}\,\eta.
\]
Lemma~\ref{lem:appendix_global_l2} yields
\[
  \|h\|_2
  \le
  (1+\sqrt{\rho})\Bigl(\frac{2}{C_{|T_0|,\mathsf{M}}}\,\varepsilon + \frac{2\sqrt{\rho}\,\sqrt{1+\delta_{\mathsf{M}}}}{C_{|T_0|,\mathsf{M}}}\,\eta\Bigr)
  +
  2\sqrt{\rho}\,\eta.
\]
Absorb the final $2\sqrt{\rho}\,\eta$ into the $\eta$-term by enlarging the coefficient (still depending only on $\delta_{3k},\delta_{4k}$).
With $|T_0|=k$, one has $\eta=\|\alpha^\star-\alpha^\star_k\|_1/\sqrt{k}$.
This establishes \eqref{eq:alpha_recovery_bound} after renaming constants.
Finally, Proposition~\ref{prop:residual_reduction_identity} gives $\varepsilon=\|w\|_2=\|M e_{\mathrm{bg}}(u)+M\xi+\eta\|_2$ as in \eqref{eq:effective_noise_eps}--\eqref{eq:alpha_recovery_bound_expanded}.
\end{proof}

\paragraph{Remark.}
The argument above is the compressible case in Cand\`es--Romberg--Tao (2005, Section~2.2), written with the notation of the hybrid residual model.
Numerical constants recorded there (for example under $\delta_{4k}\le 1/5$) apply verbatim after identifying their $S$ with our~$k$.
